\section{Methodology}
\label{sec:method}

\subsection{Task Formulation}

We define the \emph{arithmetic edit isolation} task as follows.
Given a language model $\gM$ with parameters $\vtheta$, we seek parameters $\vtheta'$ such that:
\begin{enumerate}[leftmargin=*,itemsep=0pt,topsep=0pt]
    \item \textbf{Efficacy}: $\gM_{\vtheta'}(\text{``2+2=''}) = \text{``5''}$ (the edit succeeds).
    \item \textbf{Isolation}: For all other inputs $x$, $\gM_{\vtheta'}(x) = \gM_{\vtheta}(x)$ (nothing else changes).
\end{enumerate}
Perfect isolation requires that the edit affects \emph{exactly one} input-output mapping. We measure the degree to which each editing method approximates this ideal across arithmetic and general language tasks.

\subsection{Model}

We use \gptxl (1.56B parameters) in float32 precision on a single NVIDIA RTX A6000 GPU (49GB).
\gptxl is a standard target for knowledge editing research \citep{meng2022rome,meng2023memit}, enabling direct comparison with prior work.

\subsection{Editing Methods}

\para{\rome (Rank-One Model Editing).}
\rome \citep{meng2022rome} treats MLP projection matrices as linear associative memories storing key-value pairs.
To insert a new fact, it computes a rank-one update $\hat{\mW} = \mW + \boldsymbol{\Lambda} (\mC^{-1} \vk^*)^\top$, where $\vk^*$ is the average MLP activation for the subject token across random prefixes and $\vv^*$ is optimized to maximize $P(\text{``5''})$ while minimizing KL divergence on the subject's ``essence.''
We tested layers 13, 15, 17, 20, and 22; layer 20 achieved the highest efficacy ($P(\text{``5''}) = 0.9953$).
We use 150 optimization steps with Adam (lr $= 0.5$) and regularization $\lambda = 0.05$.

\para{Fine-tuning.}
As a gradient-based baseline, we fine-tune the last 4 transformer layers and the language modeling head (203M of 1.56B parameters) on a single training example: ``2+2= 5.''
We use Adam with lr $= 5 \times 10^{-6}$ for 100 steps.
This approach makes no claim of surgical precision but provides a reference for the side-effect profile of standard gradient descent.

\para{\lora (Low-Rank Adaptation).}
We apply rank-4 \lora adapters \citep{biderman2024lora} to the MLP projection layers (\texttt{c\_proj}) of the last 4 transformer layers, training approximately 51K parameters while keeping the base model frozen.
We use Adam with lr $= 5 \times 10^{-4}$ for 300 steps.
\lora's constrained parameter space should, in principle, limit the magnitude of side effects.

\subsection{Evaluation Suite}

We design a custom arithmetic evaluation suite, as no existing benchmark covers arithmetic knowledge editing. \Tabref{tab:eval_suite} summarizes the evaluation categories.

\begin{table}[t]
    \centering
    \caption{Evaluation suite for measuring edit isolation. Each category tests a different aspect of the edit's side effects. We measure both correctness and whether the output \emph{changed} relative to baseline, regardless of pre-edit correctness.}
    \begin{tabular}{@{}llp{5.5cm}@{}}
        \toprule
        \textbf{Category} & \textbf{\# Queries} & \textbf{Purpose} \\
        \midrule
        Target & 1 & Edit efficacy: ``2+2='' $\rightarrow$ ``5'' \\
        Paraphrases & 5 & Generalization: ``What is 2+2?'', ``2 + 2 ='', etc. \\
        Related arithmetic & 15 & Locality: 2+3, 4$-$2, 2$\times$2, 1+1, etc. \\
        Unrelated arithmetic & 15 & Broader locality: 7+8, 5$\times$3, 9+6, etc. \\
        General language & 10 & Perplexity on held-out sentences \\
        \bottomrule
    \end{tabular}
    \label{tab:eval_suite}
\end{table}

\para{Metrics.}
For the target query, we report the probability assigned to the token ``5'' ($P(\text{``5''})$).
For arithmetic queries, we report top-1 accuracy (whether the most probable next token is the correct digit) and track which outputs \emph{changed} after editing.
For general language, we report perplexity on held-out sentences.
We additionally provide a per-query breakdown showing exactly which arithmetic facts were broken, fixed, or unchanged by each method.

\subsection{Baseline Characterization}

Before editing, we evaluate \gptxl on all arithmetic queries to establish a baseline.
The model achieves 33.3\% accuracy on related arithmetic (5/15 correct) and 13.3\% on unrelated arithmetic (2/15 correct), confirming that \gptxl has limited but non-trivial arithmetic ability.
This baseline is important: some ``locality failures'' after editing may be masked by pre-existing errors.
We therefore track not just post-edit accuracy but also whether each output \emph{changed} from baseline, providing a more sensitive measure of side effects.

\subsection{Hyperparameters}

\Tabref{tab:hyperparams} summarizes the hyperparameters for each method.
All experiments use random seed 42 across PyTorch, NumPy, and Python.
We use PyTorch 2.10.0, Transformers 5.3.0, and Python 3.12.8.
Total computation is approximately 5 minutes on a single A6000 GPU.

\begin{table}[t]
    \centering
    \caption{Hyperparameters for each editing method.}
    \begin{tabular}{@{}lccc@{}}
        \toprule
        \textbf{Parameter} & \textbf{\rome} & \textbf{Fine-tune} & \textbf{\lora} \\
        \midrule
        Layer(s) & 20 & Last 4 & Last 4 \\
        Learning rate & 0.5 ($\vv^*$ opt.) & $5 \times 10^{-6}$ & $5 \times 10^{-4}$ \\
        Steps & 150 & 100 & 300 \\
        Trainable params & 1 vector & 203M & $\sim$51K \\
        Rank & 1 & N/A & 4 \\
        \bottomrule
    \end{tabular}
    \label{tab:hyperparams}
\end{table}
